\chapter{Background and Related Work}
\label{ch:research}

\section{Persuasion Theory and Strategy Taxonomies}

The social-psychology literature on persuasion long predates NLP work on the
topic. Cialdini's six principles of influence -- reciprocity, commitment and
consistency, social proof, authority, liking, and scarcity
\citep{cialdini1984influence} -- are the most widely cited organizing
framework, and their fingerprints are visible in the taxonomy this project
designed (Section~\ref{sec:taxonomy}): \emph{Reciprocity and Exchange},
\emph{Commitment and Consistency}, and \emph{Urgency and Scarcity} are
category names taken almost directly from Cialdini, while \emph{Social
Influence} and \emph{Credibility Appeal} correspond to social proof and
authority respectively. The taxonomy extends this classical framework with
categories specific to charitable solicitation dialogue that do not map
cleanly onto Cialdini's six -- \emph{Rational Appeal}, \emph{Emotional
Appeal}, \emph{Framing and Presentation}, \emph{Threat and Pressure}, and
\emph{Information Provision} and \emph{Relational and Interactive} moves that
are largely conversational scaffolding rather than persuasive content per se.

\section{The Persuasion-For-Good Corpus}

\citet{wang2019persuasion} introduced the P4G corpus and, alongside it, an
11-strategy annotation scheme applied to a small labeled subset, together
with an early sequence-model strategy classifier and a persuader-side
dialogue system trained to imitate high-donation-rate human persuaders. Their
release established the corpus this project builds on, but three limitations
motivated the present work: the strategy annotation covers only a fraction of
the corpus; each turn was treated as carrying at most one dominant strategy,
whereas manual inspection of the corpus shows a turn routinely combines two
or three simultaneously (Section~\ref{sec:eda}); and no systematic,
statistically corrected analysis was published of which strategies are
actually associated with the final donation decision once dialogue-level
confounds are accounted for.

Subsequent work using P4G has largely focused on the intent-prediction side
of the corpus (predicting whether a dialogue will end in a donation) rather
than the fine-grained strategy side, which is the gap this project's
Module~B (Chapter~\ref{ch:design}) is built to fill.

\section{Multi-Label Text Classification Under Class Imbalance}

Framing strategy recognition as multi-label rather than single-label
classification is not a stylistic choice; the empirical labels---a
mean of 1.68 strategies and 1.50 categories per turn, with under 7\% of
turns carrying zero labels (Section~\ref{sec:eda})---rule out a
single-label formulation on this corpus. Binary cross-entropy per label is
the standard baseline but treats every label symmetrically, which is a poor
fit when, as here, the label distribution spans nearly three orders of
magnitude (the rarest strategy occurs twice in the corpus, the most common
occurs 2{,}860 times). \citet{benbaruch2021asl} propose Asymmetric Loss
(ASL), which decouples the focusing power applied to positive and negative
predictions and adds a small negative-probability margin so that easy
negatives contribute almost nothing to the gradient; this project adopts ASL
as its base loss (Section~\ref{sec:design-loss}) and additionally makes the
per-label positive weighting configurable per ensemble member, since it was
measured to trade micro-F1 against macro-F1 rather than improving both
(Table~\ref{tab:ablation-summary}).

Resampling is the more traditional response to label imbalance, and is what
this project's own EDA augmentation pipeline does (Section~\ref{sec:design-augmentation}).
\citet{wu2020distribution} identify a specific failure mode of resampling in
the \emph{multi-label} setting that a single-label class-balancing analysis
misses entirely: oversampling an instance to help its rare label inevitably
oversamples every \emph{other} label that instance happens to carry too, so a
turn duplicated nine times to help a rare strategy also duplicates whatever
common strategies co-occur on that turn. Their Distribution-Balanced (DB)
Loss corrects for this with a per-instance, per-class reweighting term. This
project independently measured the same failure mode on its own corpus
(Section~\ref{sec:ablation-augmentation}) -- effective rebalancing of 6$\times$
against an intended 9$\times$ -- before implementing a DB-style correction,
which is reported in the ablation study as one of the project's controlled
experiments rather than assumed to transfer from the vision literature
unmodified.

\section{Transformer Encoders for Dialogue Understanding}

The classifier ensemble (Section~\ref{sec:design-ensemble}) draws on four
pretrained encoder families: BERT-family masked language models
\citep{devlin2019bert}, RoBERTa \citep{liu2019roberta} at both base and large
scale, ELECTRA \citep{clark2020electra}, whose replaced-token-detection
pretraining objective was hypothesized to transfer well to a discriminative
turn-classification task, and TOD-BERT \citep{wu2020todbert}, which is
pretrained specifically on task-oriented dialogue and was measured to be the
one encoder that reliably converts additional dialogue context into
predictive signal on this corpus (Section~\ref{sec:ablation-context}) --
a result specific enough to the encoder's dialogue-pretraining, and not
generalizable to the other encoders, that it had to be established
empirically rather than assumed.

\section{Statistical Rigor: Conformal Prediction and Outcome Modeling}

Two further bodies of work shape the evaluation methodology
(Chapter~\ref{ch:evaluation}). Conformal prediction
\citep{vovk2005conformal} provides prediction \emph{sets} with a
distribution-free coverage guarantee rather than a single point prediction;
this project reports conformal coverage at three risk levels
(Section~\ref{sec:conformal}) as a complementary reliability check alongside
point-estimate F1. For the donation-outcome analysis, McFadden's pseudo-$R^2$
\citep{mcfadden1974conditional} is the standard goodness-of-fit measure for
logistic regression and is used throughout Section~\ref{sec:phase4-results},
always reported as a cross-validated out-of-sample estimate rather than the
optimistic in-sample figure, since with up to 28 features fitted on 1{,}017
dialogues the in-sample figure alone would overstate genuine explanatory
power.

\section{Positioning of This Work}

Relative to \citet{wang2019persuasion}, this project's contribution is an
exhaustive multi-label re-annotation of the full 10{,}600-turn corpus rather
than a sampled single-label subset, and a strategy-to-outcome analysis run on
that exhaustive annotation with cross-validated, multiple-testing-corrected
statistics rather than raw association tests. Relative to the multi-label
classification literature \citep{benbaruch2021asl, wu2020distribution}, the
contribution is empirical rather than methodological: existing techniques
(ASL, DB-style rebalancing, empirical-Bayes threshold shrinkage) are applied
and, critically, \emph{measured} on this specific corpus rather than assumed
to transfer, and the report documents several cases (Chapter~\ref{ch:evaluation})
where a technique that is well established elsewhere measurably failed to
help, or helped one metric while hurting another, on this dataset --
findings that are only visible because of the controlled, sentinel-verified
experimental protocol described in Section~\ref{sec:design-methodology}.
